\documentclass[12pt,a4paper]{article}

\usepackage[utf8]{inputenc}
\usepackage[T1]{fontenc}
\usepackage{lmodern}
\usepackage{microtype}
\usepackage{amsmath,amssymb}
\usepackage[colorlinks=true,linkcolor=black,citecolor=black,urlcolor=black]{hyperref}
\usepackage[numbers,sort&compress]{natbib}
\usepackage{geometry}
\usepackage{parskip}
\usepackage{setspace}
\usepackage{titlesec}
\usepackage{booktabs}

\geometry{margin=2.5cm}
\setstretch{1.15}

\titleformat{\paragraph}[runin]{\normalfont\normalsize\bfseries}{}{0em}{}[.]
\titlespacing{\paragraph}{0pt}{0.75\baselineskip}{0.5em}

\title{\textbf{A Cooperative Signalling-Game Formulation for\\[0.3em]
Multi-Agent Highway Driving}\\[0.6em]
\large Measuring Emergent Behavioural Conventions via\\
Multi-Agent Reinforcement Learning}
\author{Nhan Nguyen}
\date{Master's Dissertation Proposal}

\begin{document}

\maketitle

\begin{abstract}
Autonomous highway driving is, at its core, a multi-agent decision problem: every vehicle
carries a private behavioural type that can't be directly observed. Two broad approaches have
dominated the field. The first infers types from observed trajectories but doesn't account for the
fact that the ego-vehicle is simultaneously being inferred by others. The second solves analytical
equilibria under assumed payoffs with hand-specified signals. Neither asks what the classical
signalling-game literature treats as primary - given a shared behavioural signal space, what
communicative conventions actually emerge from interaction?

This thesis takes that question seriously. Cooperative multi-agent highway driving is formulated
as a repeated signalling game over a small discrete vocabulary of behavioural signals whose
meanings aren't pre-assigned, and the methodological machinery is developed to measure
emergent informativeness via multi-agent reinforcement learning with a recurrent type encoder. Three concrete deliverables follow: (i) a formal cooperative signalling-game model for
continuous-control highway driving with a behaviourally-grounded, non-cheap-talk signal space;
(ii) an operational mutual-information criterion, $\rho$, for detecting pooling, partial-pooling, and
separating regimes; and (iii) an experimental demonstration across three conditions establishing
that $\rho$ responds to the behavioural content of type labels rather than their mere presence,
and that the predominant pooling at the current preference divergence ($d=0.30$) is consistent
with the Spence/Crawford--Sobel prediction that separation requires $d > d_{\text{sep}}$, a
threshold derived analytically from the reward structure's single-crossing condition.
\end{abstract}

\section{Introduction}

Highway driving means constant negotiation with vehicles you can't fully read. Caution,
assertiveness, risk tolerance, preferred following distance - none of it is directly observable.
The field's standard answer has been perception, prediction, planning: build a probabilistic model
of neighbour motion and let the planner marginalise over it~\cite{chen2022milestones,yurtsever2020survey}.
Prediction methods have come a long way - from rule-based forward simulation, through belief
tracking over discrete intent hypotheses~\cite{bai2015intention,zhou2018joint}, to unsupervised
latent-trait inference with recurrent or attention-based
encoders~\cite{liu2022learning,liu2024automatic,ma2021reinforcement}.

But that entire trajectory frames the problem one-sidedly. The ego-vehicle asks \textit{what
will my neighbour do?} and stops there. A human driver also asks the other half:
\textit{what do my neighbours think I am, and is my behaviour shaping that belief?}
Real drivers don't just execute manoeuvres at minimum cost. The way they close a gap, ease
off the accelerator, or begin a lane change typically carries more information than what's strictly
needed to complete the move. Those excess margins are signals, and other
drivers read them~\cite{fox2019chicken,risto2017human,wang2020driver}.

Consider a concrete two-vehicle instance. Car~A is cautious; car~B is assertive. Both
types are private: neither vehicle can directly observe the other's behavioural disposition.

B approaches A from behind in the same lane. Over the preceding five seconds, B has
maintained a tight following distance, accepted small gaps, and initiated accelerations
sharply. A observes this kinematic history and updates its belief: the follower is probably
assertive. A responds by holding a wider headway than it otherwise would, smoothing its
own deceleration profile to avoid provoking a rear-end conflict. Critically, A did not change
its \textit{destination} - it is still tracking the same speed target - but it changed
\textit{how} it pursues that target, because its model of B has changed.

Meanwhile, B is reading A. A's wide margins, early braking, and smooth acceleration
smoothness are recognised as the kinematic signature of a cautious type. B updates its
belief accordingly and adjusts its gap-acceptance threshold upward: it will wait for a
larger gap before forcing a merge, because it has learned that A will not accelerate
assertively to close one. Again, B's destination is unchanged; only its execution style
shifts in response to its inferred model of A.

Neither vehicle has sent an explicit message. Coordination emerged entirely from the
observable style of otherwise routine manoeuvres. The same physical action - maintain
speed, decelerate, drift laterally - carried different information depending on
\textit{how} it was executed, and both vehicles exploited that information to reduce
conflict without any direct communication channel.

This is the province of signalling-game theory~\cite{crawford1982strategic,spence1973job}.
Informed senders transmit observable signals; uninformed receivers update beliefs via Bayes'
rule and respond accordingly. Equilibria fall into three canonical regimes: \textit{separating},
where different types emit different signals and communication is fully informative;
\textit{pooling}, where all types converge on one signal and communication collapses; and
\textit{partial-pooling} somewhere in between~\cite{fudenberg1991game}. The existence of a
signal space doesn't determine by itself whether it'll be used informatively. Whether it is depends
on the game structure.

This thesis casts cooperative multi-agent highway driving as such a signalling game, then builds
the empirical tools to determine which regime actually emerges when agents are trained over a
shared discrete vocabulary whose meanings are deliberately left unassigned. The
cooperative-population restriction is methodologically deliberate: it removes strategic deception
as a confound. Any pooling outcome can't be attributed to agents intentionally misrepresenting
themselves, so the presence or absence of separation becomes diagnostic of game structure rather
than adversarial strategy.

% \section{Research Gap}
%
% The gap this work targets is visible across three parallel literatures whose combined coverage
% still doesn't reach our question.
%
% \paragraph{Strand 1: Analytical signalling games in lane-changing}
% \citet{shao2020discretionary} model discretionary lane change as a signalling game between
% aggressive and conservative sender types, calibrating a Perfect Bayesian Equilibrium against
% NGSIM data. The signal is a single binary indicator; payoffs are hand-specified; the equilibrium
% is solved, not learned. A broader analytical
% literature~\cite{ali2019game,talebpour2015modeling,zhang2022game} likewise assumes fixed
% payoffs and solves for Nash or Stackelberg equilibria. This establishes theoretical precedent but
% leaves open how a signalling policy would be acquired from experience rather than imposed -
% and what happens when the signal space is larger than binary with semantics that aren't
% pre-assigned.
%
% \paragraph{Strand 2: Latent-trait inference plus RL planning}
% A growing body of work infers driver traits via variational or attention-based encoders and
% concatenates the latent representation with the planning state in an RL
% policy~\cite{chakraborty2023structural,liu2022learning,liu2024automatic,ma2021reinforcement}.
% These methods are structurally one-sided: inference flows from neighbours to the ego-vehicle
% only. There's no construct of the ego-vehicle projecting a type outward, no signal space separate
% from the action space, no equilibrium concept at all.
%
% \paragraph{Strand 3: Emergent communication in MARL}
% The MARL literature on emergent
% communication~\cite{foerster2016learning,lazaridou2020emergent,sukhbaatar2016learning} asks
% precisely the question we raise - how agents learn to use a shared signal space. But it almost
% exclusively assumes a cheap-talk channel: messages carry no physical cost. Driving violates this.
% Every signal must be carried on physical behaviour that simultaneously affects the trajectory and
% incurs driving costs. Costly signalling~\cite{spence1973job,zahavi1975mate} produces
% fundamentally different equilibrium structures than cheap talk does.

% \paragraph{Statement of the gap}
% Strand~1 gives the equilibrium framework but no learning. Strand~2 gives learning-based type
% inference but no signalling. Strand~3 gives learned signalling but only under cheap-talk
% assumptions that don't hold for driving. What's missing is a formulation that is simultaneously:
% (a) a signalling game in the Spence/Crawford--Sobel sense; (b) learned end-to-end via
% decentralised MARL; (c) played over a discrete, behaviourally-grounded vocabulary with
% unassigned semantics; and (d) studied for the emergent convention itself, not just for the
% resulting driving policy.

\section{Research Questions}

\paragraph{RQ1 - Measurable emergence of non-trivial signal structure}
When cooperative agents are trained jointly with a shared discrete behavioural signal space over
a heterogeneous population, does the separation coefficient $\rho$ (Section~\ref{sec:eval})
reliably exceed a random-signal baseline across training seeds?

\paragraph{RQ2 - Necessity of type heterogeneity}
Does $\rho$ differ between a bimodal population where type labels correspond to distinct reward
functions and one where both labels share identical reward weights? In the same-reward control,
both type tokens are present and sampled equally so $\rho$ remains well-defined, but the labels
carry no behavioural content - no incentive gradient exists to drive separation. If $\rho$ stays
near the random-signal level in this control but rises in the bimodal condition, any elevation is
attributable to the behavioural content of the type labels rather than to architectural effects of
the signalling head.

\section{Method}

\subsection{Formalism}

The highway is modelled as a partially observable stochastic game. Each agent $i$ has a latent
type $\theta_i \in \Theta$, fixed within an episode and resampled across episodes from a
population distribution. We use $|\Theta| = 2$ (cautious, assertive), with the cluster
determining the weighting of safety, progress, and comfort terms in the agent's intrinsic reward.
Types are private; there are no persistent identifiers across episodes, and no communication
channel exists other than the behavioural signal space defined below.

At each decision step $t$, agent $i$ selects a driving action $a_i \in \mathcal{A}$ -
continuous: longitudinal acceleration, steering rate - and a discrete signal $m_i \in \mathcal{M}$
from a vocabulary of cardinality $|\mathcal{M}| = 3$. The signal is a discrete selector that
determines which of three pre-specified execution profiles renders $a_i$. Every profile is
physically feasible and driving-legal; what differs across them is the observable kinematic
signature.

A joint policy is a pair $(\sigma_i, \pi_i)$ for each agent: a signalling strategy $\sigma_i(s_i,
\theta_i) \to \Delta(\mathcal{M})$ and a response strategy $\pi_i(s_i, b_i) \to \Delta(\mathcal{A})$,
where $b_i$ is the belief held by agent $i$ over the types of currently visible neighbours. The
analytical reference point is Perfect Bayesian Equilibrium~\cite{fudenberg1991game}: strategies
are mutually best-responding, and on-path beliefs are Bayes-updated as
\begin{equation}
    b_i(\theta_j \mid m_j, h_j) \;\propto\; \sigma_j(m_j \mid s_j, \theta_j)\,
    b_i(\theta_j \mid h_j),
    \label{eq:bayes}
\end{equation}
where $h_j$ is the recent observed behavioural history of agent $j$. We don't claim to find a
PBE; we train via MARL and characterise whatever equilibrium the training procedure reaches.

\subsection{Architecture}

For each visible neighbour $j$, a type embedding is produced by a single-layer GRU encoder
over a windowed sequence of observed kinematic features and emitted signals:
\[
    \hat{\theta}_{ij}(t) \;=\; f_\phi\!\left(\bigl\{(o_j(s),\, m_j(s))\bigr\}_{s=t-W}^{t}\right),
\]
with window $W = 5\,\text{s}$. A GRU is chosen as the primary encoder for training stability at
this scale. The encoder is shared across agents under parameter sharing during centralised training; at execution time every agent runs its own copy over its own observations.

The full policy decomposes into three heads sharing the encoder. There's the type encoder
$f_\phi$; a signalling head $\sigma_\theta$ mapping $(s_i, \theta_i, \{\hat{\theta}_{ij}\}_j)$ to
a categorical distribution over $\mathcal{M}$ on the sender side; and a response head $\pi_\psi$
mapping $(s_i, \{\hat{\theta}_{ij}\}_j)$ to a continuous driving action on the receiver side - this
head doesn't take the agent's own type as input. Neighbour embeddings are aggregated by a
permutation-invariant DeepSets layer~\cite{zaheer2017deep} to handle variable neighbour counts.
Training uses MAPPO~\cite{yu2022surprising} under centralised training with decentralised
execution~\cite{lowe2017multi}. MAPPO is the natural choice here: the signalling head is
categorical, and MAPPO is the established default for cooperative MARL with mixed
discrete-continuous action spaces. The indirect credit-assignment problem for the signalling
head --- gradient arrives only through downstream neighbour behaviour --- is mitigated by two
mechanisms: the centralised critic provides a variance-reducing value baseline for the
REINFORCE gradient through the discrete signal, and a separate entropy coefficient on the
signalling head is annealed linearly from a high initial value to a low final value over
training, maintaining signal exploration until type-informative structure can form.

\subsection{Signal Space}

$\mathcal{M}$ is a small discrete vocabulary of fixed execution profiles. For any commanded
driving action $a$, each $m \in \mathcal{M}$ specifies a distinct but physically feasible way to
realise $a$. Each profile is a parametric override passed to the IDM/MOBIL low-level controllers
in \texttt{highway-env}: the commanded action $a_i$ sets the target, and the profile determines
the kinematic envelope within which it is executed. The three profiles are defined by the
parameters in Table~\ref{tab:profiles}.

\begin{table}[h]
\centering
\small
\begin{tabular}{lcccc}
\toprule
Profile & Headway & Jerk cap & Gap-accept & Lateral drift \\
        & multiplier & (m\,s$^{-3}$) & threshold (m) & at initiation (m) \\
\midrule
$m_1$ (conservative) & $1.5\times$ & 1.0 & 15 & 0.3 \\
$m_2$ (neutral)      & $1.0\times$ & 2.0 & 10 & 0.6 \\
$m_3$ (assertive)    & $0.7\times$ & 4.0 &  6 & 1.0 \\
\bottomrule
\end{tabular}
\caption{Execution-profile parameters. Headway multiplier scales the IDM desired headway;
jerk cap bounds longitudinal acceleration rate; gap-accept threshold is the minimum accepted
merge gap; lateral drift is the anticipatory displacement at lane-change initiation.}
\label{tab:profiles}
\end{table}

Two design commitments follow from this. First, signal emission is costly in the sense of
\citet{spence1973job}: different profiles incur different costs across progress, comfort, and
safety dimensions, and no single profile dominates for every type. This cost asymmetry is what
gives a separating equilibrium its leverage - cautious types find it cheaper to emit cautious
profiles, assertive types find it cheaper to emit assertive ones - and it's exactly what a
cheap-talk channel lacks. Second, and just as importantly, the semantics of each signal are
deliberately left unassigned. No profile is labelled as meaning \textit{I am cautious}. Whether
$m_k$ comes to denote a particular type is an emergent property of the trained population,
measured rather than stipulated.

\subsection{Reward Structure}

The reward for agent $i$ is $R_i = R^{\mathrm{drive}}_i$, the standard driving reward over
forward progress, collision avoidance, lane-keeping, and comfort, weighted according to
$\theta_i$~\cite{leurent2018environment}. There's no reward term for the signal: no explicit
bonus for truthful signalling, no penalty for lying. This is a deliberate methodological
commitment. If we rewarded truthful signalling directly, RQ1 would be answered by
construction - separation would emerge because we paid for it. Keeping $R_i$ free of any
signal-dependent term ensures the regime we observe is an endogenous consequence of
interaction structure, not reward engineering.

This means the signalling head receives gradient only indirectly, through the effect of its signal
on neighbour behaviour, which in turn affects $R^{\mathrm{drive}}_i$. It's the same
credit-assignment pattern that emergent-communication work~\cite{foerster2016learning} must
solve - but with the added complexity that our messages also have direct physical consequences
for the sender. We expect this credit-assignment difficulty to be the principal practical challenge
of the approach.

\subsection{Preference Divergence and Separating Equilibrium Threshold}

The degree to which the two types face misaligned incentives is captured by a scalar
preference divergence:
\begin{equation}
    d \;=\; 1 - \frac{w_c \cdot w_a}{\|w_c\|\,\|w_a\|},
    \label{eq:divergence}
\end{equation}
where $w_\theta = (w^\text{safety}_\theta,\, w^\text{progress}_\theta,\,
w^\text{comfort}_\theta)$ is the reward-weight vector for type $\theta$. $d = 0$ when both
types share identical weights (the same-reward control); $d$ increases as the weight vectors
diverge.

The Crawford--Sobel / Spence single-crossing condition for a separating equilibrium requires
that the reward differential of using profile $m_1$ versus $m_3$ has opposite signs across
types. Define the profile differential vector
$\boldsymbol{\Delta} = (\Delta_s, \Delta_p, \Delta_k)$, where
$\Delta_i = \mathbb{E}[r_i \mid m_1] - \mathbb{E}[r_i \mid m_3]$ is the expected per-step
reward component under profile $m_1$ minus that under $m_3$. From Table~\ref{tab:profiles}:
$\Delta_s > 0$ (m$_1$ is safer), $\Delta_p < 0$ (m$_1$ is slower), and $\Delta_k > 0$
(m$_1$ is smoother). Separation is supportable iff
\begin{equation}
    w_c \cdot \boldsymbol{\Delta} > 0
    \quad \text{and} \quad
    w_a \cdot \boldsymbol{\Delta} < 0.
    \label{eq:single-crossing}
\end{equation}

To derive the minimum divergence $d_\text{sep}$ at which~\eqref{eq:single-crossing} first
holds simultaneously, parameterise the weight vectors as
\begin{equation}
    w_c(\alpha) = \bar{w} + \alpha\,\delta, \qquad
    w_a(\alpha) = \bar{w} - \alpha\,\delta,
    \label{eq:weight-param}
\end{equation}
where $\bar{w} = \tfrac{1}{2}(w_c + w_a)$ is the mean weight vector and
$\delta = \tfrac{1}{2}(w_c - w_a)$ is the half-difference; $\alpha = 0$ gives the
same-reward baseline and $\alpha = 1$ recovers the current bimodal configuration.
Each inequality in~\eqref{eq:single-crossing} imposes a linear threshold on $\alpha$:
\begin{equation}
    \alpha_c \;=\; \frac{-\bar{w} \cdot \boldsymbol{\Delta}}{\delta \cdot \boldsymbol{\Delta}},
    \qquad
    \alpha_a \;=\; \frac{\;\bar{w} \cdot \boldsymbol{\Delta}}{\delta \cdot \boldsymbol{\Delta}},
    \qquad
    \alpha_\text{sep} \;=\; \max(\alpha_c,\, \alpha_a).
    \label{eq:alpha-sep}
\end{equation}
Converting back to cosine distance gives
$d_\text{sep} = 1 - \cos\!\bigl(w_c(\alpha_\text{sep}),\, w_a(\alpha_\text{sep})\bigr)$.
If the current bimodal condition satisfies $d > d_\text{sep}$, the formulation lies in the
analytically supportable separation regime. The $d$-sweep in Section~\ref{sec:conditions} tests whether empirical $\rho$ tracks this
boundary across five divergence levels. Note that $d$ and $d_{\text{sep}}$ are our
operationalisations of the preference-bias concept introduced by~\citet{crawford1982strategic};
the single-crossing condition in~\eqref{eq:single-crossing} is the discrete analogue of
Spence's~\cite{spence1973job} costly-signalling separability criterion applied to our reward
structure.

\subsection{Environment}

The environment is \texttt{highway-env}~\cite{leurent2018environment}, extended to support
(i)~heterogeneous agent populations drawn from a configurable type distribution, and (ii)~the
discrete execution-profile vocabulary implemented as low-level controllers layered on the
existing continuous action space. A single scenario is used throughout: dense multi-lane highway
driving with $N = 6$ agents at nominal traffic density.

\section{Evaluation Protocol}
\label{sec:eval}

\subsection{Primary Diagnostic}

The central measurement is the empirical mutual information between the latent type and the
emitted signal, estimated from evaluation rollouts:
\begin{equation}
    I(\Theta;\, \mathcal{M}) \;=\; \sum_{\theta \in \Theta,\, m \in \mathcal{M}}
    p(\theta, m)\log\frac{p(\theta, m)}{p(\theta)\,p(m)}.
    \label{eq:mi}
\end{equation}
where $\Theta = \{\text{cautious},\,\text{assertive}\}$ is the type space; $\mathcal{M} = \{m_1, m_2, m_3\}$ is the signal vocabulary; $p(\theta, m)$ is the empirical joint frequency of type $\theta$ emitting signal $m$, estimated from evaluation rollouts; $p(\theta)$ is the marginal type frequency (fixed at $0.5$ each in the bimodal condition); $p(m)$ is the marginal signal frequency, aggregated across all types; and the log-ratio $\log \tfrac{p(\theta,m)}{p(\theta)\,p(m)}$ measures how much more often the pair $(\theta, m)$ co-occurs than expected under statistical independence - zero when type and signal are unrelated, positive when a type preferentially selects a signal.

We report $I(\Theta; \mathcal{M})$ in bits and normalise by $\log|\Theta|$ to give a separation
coefficient $\rho \in [0, 1]$. Operationally: $\rho < 0.1$ indicates pooling (the signal carries
essentially no information about type); $\rho > 0.9$ indicates separating (the signal determines
type); intermediate values indicate partial-pooling. The distribution of $\rho$ across seeds and
conditions directly addresses RQ1 and RQ2.

Per-signal conditional distributions $p(\theta \mid m)$ are reported as a supplementary
diagnostic, revealing whether any partial-pooling structure is clean - some signals fully
informative, others uninformative - or gradual, where all signals carry intermediate information.

\subsection{Intervention Test}
\label{sec:intervention}

Mutual information shows that signals are correlated with type but does not establish that
receivers exploit those signals. A correlation - cautious agents happen to emit $m_1$
for kinematic reasons unrelated to any communicative act - could produce elevated $\rho$ with
no actual signalling function. To test causal use, the following seven conditions are evaluated
on the best bimodal checkpoint, each over 200 episodes:

\begin{enumerate}
    \item \textbf{Intact (baseline).} Unmodified policy; all ablations are compared against
    this reference on collision rate ($\Delta_\text{collision}$) and merge-success rate
    ($\Delta_\text{merge}$).

    \item \textbf{Signal randomisation.} Each emitted token is replaced with a uniform draw
    from $\mathcal{M}$ before being passed to the environment and to neighbour observations.
    If receivers condition on token content, policy quality should degrade. $\rho$ collapses by
    construction here; the relevant measure is driving-performance degradation.

    \item \textbf{Signal hiding.} The signal component is zeroed out in each neighbour's
    observation vector, so receivers see kinematics only with no token. This tests whether the
    response head uses the signal channel at all.

    \item \textbf{Signal permutation.} The mapping from token to execution profile is scrambled,
    so the kinematic consequence of each token no longer matches what was learned. This separates
    the communicative token from its physical referent.

    \item \textbf{Kinematic hiding.} Kinematic observations are zeroed while signal tokens are
    left intact. Receivers see tokens only, with no co-occurring kinematics. This tests whether
    tokens carry information independently of the kinematic context in which they are emitted.

    \item \textbf{Full blackout.} Both signal tokens and kinematic observations are zeroed.
    This establishes floor performance when receivers are entirely blind, providing a lower bound
    against which the partial hiding conditions are interpreted.

    \item \textbf{Desync.} The policy token is emitted and remains visible in neighbour
    observations as normal, but a uniformly random execution profile is applied kinematically
    at each step via a low-level override. This breaks the token$\leftrightarrow$kinematic-image
    binding without altering what receivers observe in the token channel. A crash spike under
    desync - but not under signal hiding - indicates that receivers relied on the kinematic image
    of the token rather than the token itself; the absence of such a spike indicates that the token
    channel is the primary carrier. To our knowledge, no prior emergent-communication study
    includes this ablation: cheap-talk architectures have no separable physical-execution
    component, so the token--kinematic decoupling that desync exploits cannot arise.
\end{enumerate}

The critical diagnostic is the comparison between conditions~3 and~7: elevated $\rho$ that
survives signal hiding ($\rho_\text{hide} \approx \rho_\text{intact}$) indicates the
type--signal correlation is kinematic rather than communicative; a desync crash spike confirms
that receivers had been exploiting the kinematic image of the token. Conditions~3--6 jointly
decompose the signal channel into its token and kinematic components and establish whether each
independently supports coordination.

\subsection{Experimental Conditions}
\label{sec:conditions}

Three conditions are compared, each across twenty training seeds:
\begin{enumerate}
    \item \textbf{Full method, bimodal population.} 50\% cautious, 50\% assertive. The
    principal condition; addresses RQ1.

    \item \textbf{Full method, same-reward bimodal population.} Both type labels \{cautious,
    assertive\} are present and sampled equally ($50\%$ each), but both receive the same reward
    weight vector - identical safety, progress, and comfort coefficients. The signalling head
    observes a type token, but the token carries no behavioural content. Tests RQ2: $\rho$ should
    remain near the random-signal level here but be elevated in condition~1 if type-behavioural
    content, not the label itself, drives separation.

    \item \textbf{Random-signal baseline, bimodal population.} Identical architecture with the
    signalling head replaced by uniform random emission over $\mathcal{M}$. Controls for the
    possibility that any measured $\rho$ is an artefact of the extra action dimension rather than
    learned informativeness.
\end{enumerate}

\subsection{Statistical Protocol}

Results are reported as per-seed values with means and 95\% confidence intervals across twenty seeds. Evaluation
uses 200~episodes per seed. Driving-performance metrics - collision rate, near-miss count,
merge-success rate - are logged throughout for contextual reporting.

% \subsection{Toy-Configuration Sanity Check}
%
% Before committing compute to the three main conditions, a two-type / two-signal toy
% configuration ($|\Theta| = 2$, $|\mathcal{M}| = 2$) is trained to verify that $\rho$ rises above
% the random-signal baseline over the course of training. This serves as an early architectural
% validation: if $\rho$ doesn't move on the toy configuration, the system needs diagnosis before
% full-scale runs are attempted.
%
% \section{Theoretical Interpretation}
%
% No formal PBE existence proof is attempted for the full game. In continuous control this isn't
% tractable, and even if it were, it wouldn't tell us which equilibrium a learning procedure selects
% among potentially many. Instead, the empirical findings are related to the analytical
% signalling-game literature in one specific way.
%
% \citet{crawford1982strategic} show that in a one-shot sender--receiver game, the number of
% distinguishable messages at any equilibrium is bounded above by a quantity that decreases as
% sender and receiver preferences diverge. In the perfectly aligned case, arbitrarily fine
% communication is possible; in the perfectly misaligned case, only pooling equilibria survive. In
% our cooperative setting, all agents share the same driving reward structure up to type-dependent
% weights - preferences are mostly aligned but not identically so, since a cautious sender and an
% assertive receiver still disagree about the value of a tight gap. Crawford--Sobel's qualitative
% prediction is therefore that partial-pooling, rather than full separation, should be the generic
% outcome. That coincides with H1.
%
% We report $\rho$ and the conditional distributions $p(\theta \mid m)$ for each condition and
% state whether the observed values are qualitatively consistent with this prediction. The goal is a
% grounded, non-overclaiming connection between the measured emergent regimes and the
% classical theory.

\section{Expected Contributions}

\begin{enumerate}
    \item \textbf{Formulation.} A cooperative repeated signalling-game formulation of
    multi-agent highway driving with a behaviourally-grounded, non-cheap-talk signal space,
    closing the gap between analytical signalling-game lane-change
    models~\cite{shao2020discretionary} and learning-based MARL.

    \item \textbf{Operational diagnostic.} The separation coefficient $\rho$, grounded in
    empirical mutual information, as a method-independent criterion for classifying trained MARL
    populations as pooling, partial-pooling, or separating; complemented by a seven-battery
    intervention test --- including a novel desync condition that keeps the emitted token visible
    to receivers while randomising the sender's physical execution, an ablation that cannot be
    constructed in cheap-talk architectures --- that moves beyond correlation to causal use of
    the signal channel.

    \item \textbf{Experimental demonstration.} A three-condition experiment establishing that
    the formulation trains end-to-end, that $\rho$ responds to the behavioural content of type
    labels rather than their mere presence, that the predominant pooling regime at the current preference divergence ($d=0.30$, mean
    $\rho=0.087$) is consistent with the analytically derived threshold $d_{\text{sep}}$ ---
    separation is not expected below this value --- and that $\rho$ is elevated 43-fold above
    the same-reward control across all seeds, attributing measured separation to the
    behavioural content of type labels rather than to architectural or labelling effects.
    A transformer encoder is identified as a future extension if the GRU baseline trains reliably.
\end{enumerate}

\section{Practical Relevance}

Informal behavioural signalling is part of how human drivers coordinate - and the absence of a
modelled signalling layer is one plausible source of the stilted or freezing interactions sometimes
seen in deployed autonomous systems. A driving stack that can emit legible behavioural cues and
read them is likely to interact more fluidly with human traffic, regardless of whether those
specific signals emerge from training or are deliberately designed. This thesis provides the
measurement apparatus needed to ask whether a given MARL driving policy has learned a
legible behavioural vocabulary at all - which is a prerequisite for deliberate design.


\bibliographystyle{plainnat}
\bibliography{references}

\end{document}
